\documentclass[11pt]{article}
\usepackage[margin=2.4cm]{geometry}
\usepackage{amsmath,amssymb,amsthm,booktabs,graphicx,xcolor,hyperref}
\newcommand{\bb}{\boldsymbol\beta}
\newcommand{\X}{\mathbf{X}}
\newcommand{\y}{\mathbf{y}}
\newcommand{\I}{\mathbf{I}}
\newcommand{\B}{\mathbf{B}}
\newcommand{\Lm}{\mathbf{L}}
\newcommand{\Dm}{\mathbf{D}}
\newcommand{\G}{\mathbf{G}}
\newcommand{\bv}{\mathbf{b}}
\newcommand{\tr}{\operatorname{tr}}
\newcommand{\fig}[2]{\IfFileExists{#1}{\includegraphics[width=#2\linewidth]{#1}}{\fbox{\parbox{0.9\linewidth}{\centering\small\texttt{#1}}}}}
\newtheorem{remark}{Remark}

\title{Bayesian rows for the modified-Cholesky EnKF:\\ climatological priors and adaptive memory in the precision estimate}
\author{}
\date{Draft, 15 September 2026}

\begin{document}\maketitle

\begin{abstract}
The modified-Cholesky ensemble Kalman filter (EnKF-MC) estimates the background
precision $\B^{-1}=\Lm^\top\Dm\Lm$ one row at a time, each row being a regression of a
state component on its predecessors within a localization radius, fitted to the $N$
ensemble anomalies with a ridge penalty. We treat each row as a Bayesian regression and
change two things about its prior. First, the prior mean is not zero: at the start of a
cycled experiment it is the coefficient vector of the same regression fitted to a
climatology of the model, and from the second cycle on it is the previous cycle's
posterior. Second, the prior weights are not fixed: the forgetting factor of the
sequential posterior and the weight of the climatological term are chosen for every row
and every cycle by maximising the marginal likelihood of the new ensemble anomalies.
The result is a precision whose coefficients are estimated from an effective sample far
larger than $N$ where the forecast error is stationary, and from the current ensemble
alone where it changes. On the quasi-geostrophic testbed of Sakov and Oke with $N=40$
members, the Bayesian rows reduce the analysis error of the plain EnKF-MC at every
cycle of a 36-cycle run (by 5--20\%), reduce the error of the first analysis from a
climatological ensemble by 6--9\% at every observation density between 1\% and 40\%,
and end four times below the best LETKF. The hyperparameters selected by the evidence
have a transparent physical reading: the climatology is used over the jet and released
after one analysis; memory is kept where the flow is quiescent and dropped along the
jet. We also report, with the same care, what does not work: five ways of choosing
\emph{which} predecessors a row uses, from the ensemble, from the dynamics or from the
climatology, none of which beats a uniform radius once the filter has converged; the
gain is in the values of the rows and in how much data they see, not in their sparsity
pattern.
\end{abstract}

\section{Introduction}

Localization in ensemble Kalman filters exists because $N$ members cannot estimate an
$n\times n$ covariance. The modified-Cholesky EnKF (EnKF-MC) of Nino-Ruiz and
co-workers localizes in the precision: component $i$ is regressed on its predecessors
$P(i)$, the components with a lower index within a radius $r$, and the regression
coefficients fill row $i$ of a unit lower-triangular $\Lm$. With a ridge penalty, the
row is a well-posed problem even when $|P(i)|>N$, the precision is sparse, and the
analysis is a sparse linear solve. The radius is the localization parameter; it is
chosen by a sweep, and Bickel and Levina's theory of banded Cholesky estimators says
why a small radius is right when $N$ is small: the bandwidth balances the bias of
truncating the true dependence against the variance of estimating it from $N$ samples.

This paper does not change the sparsity pattern. It changes what the rows are fitted
to. Two observations motivate it. The first is that the ridge penalty of a row is a
Gaussian prior on its coefficients, centred at zero; nothing forces the centre to be
zero, and at the first analysis of a cycled experiment, when the ensemble is a sample
of the model's climatology, the coefficients of the same regression fitted to that
climatology are the correct prior mean. The second is that the forecast error at one
cycle resembles the forecast error at the previous one, so the previous row is a
useful prior for the current one; how useful depends on how fast the error is changing,
which differs from place to place and can be read from the data. Both observations
lead to the same construction: a sequential Bayesian regression per row, whose
hyperparameters -- a forgetting factor and a climatological weight -- are chosen by
the marginal likelihood of the new anomalies.

Section~\ref{sec:method} gives the method. Section~\ref{sec:setup} the testbed.
Section~\ref{sec:results} the results: first analysis, cycled runs, comparison with
the LETKF, and the maps of what the evidence decides. Section~\ref{sec:negative}
reports the structural alternatives that were tried and why they lose.
Section~\ref{sec:discussion} discusses limits and what is left to establish.

\section{Method}\label{sec:method}

\subsection{The modified-Cholesky row as a regression}
Let $\Delta\X\in\mathbb{R}^{n\times N}$ be the forecast anomalies. For component $i$
with predecessors $P(i)$ ($p=|P(i)|$), write $\y\in\mathbb{R}^N$ for its anomalies and
$\X\in\mathbb{R}^{N\times p}$ for those of its predecessors. The row model is
\begin{equation}
  \y = \X\bb + \boldsymbol\varepsilon,\qquad \boldsymbol\varepsilon\sim\mathcal N(0,\sigma^2\I),
\end{equation}
and $\Lm_{ii}=1$, $\Lm_{i,P(i)}=-\bb^\top$, $\Dm_{ii}=1/\sigma^2$, $\widehat\B^{-1}=\Lm^\top\Dm\Lm$.
The standard estimator is ridge regression,
\begin{equation}
  \widehat\bb = (\X^\top\X + a\I)^{-1}\X^\top\y,\qquad a=\alpha\,\tr(\X^\top\X)/p,
  \label{eq:ridge}
\end{equation}
with $\alpha$ dimensionless and the same for every row. This is the posterior mean of
$\bb$ under the prior $\bb\sim\mathcal N(0,\sigma^2/a\,\I)$.

\subsection{Bayesian rows}
We keep the likelihood and change the prior. All statistics are kept in correlation
scale -- $\X$ and $\y$ standardised by their ensemble standard deviations -- so that
cycles with different spread contribute equally; coefficients are rescaled at the end.

\paragraph{Sequential prior.} Let $\G_t=\X_t^\top\X_t/N$ and $\bv_t=\X_t^\top\y_t/N$
be the sufficient statistics of cycle $t$. The row keeps a memory
\begin{equation}
  \G_t^{\rm mem} = \rho_t\,\G_{t-1}^{\rm mem} + \G_t,\qquad
  \bv_t^{\rm mem} = \rho_t\,\bv_{t-1}^{\rm mem} + \bv_t,
  \label{eq:memory}
\end{equation}
with forgetting factor $\rho_t\in[0,1)$: the posterior of one cycle is the prior of the
next, discounted.

\paragraph{Climatological term.} From $M$ independent states of the model's climatology
the same statistics are computed once, $\G^{c}$, $\bv^{c}$, weighted by $w_0=M/N$: the
climatology is worth as many cycles as ensembles fit in it. At cycle $t$ the row's prior
is $\G^{\rm mem}_t + w_t\G^c$, $\bv^{\rm mem}_t+w_t\bv^c$ with $w_t\in[0,w_0]$.

\paragraph{Estimator.} The row coefficients are
\begin{equation}
  \widehat\bb_t = \big(\G^{\rm mem}_t + w_t\G^c + a\,\mathbf{T}\big)^{-1}\big(\bv^{\rm mem}_t+w_t\bv^c\big),
  \qquad \mathbf T=\operatorname{diag}\big((1+\tau d_{ij})^2\big),
  \label{eq:bayesrow}
\end{equation}
where $d_{ij}$ is the Chebyshev distance from $i$ to predecessor $j$ and $\tau$ a
tapering slope: the penalty of a coefficient grows with the distance of its predecessor,
a soft version of the banded Cholesky. The residual variance $\sigma^2$ is that of the
current ensemble under $\widehat\bb_t$, so the scale of $\widehat\B^{-1}$ follows the
ensemble while its shape draws on the memory and the climatology.

\paragraph{Hyperparameters by evidence.} $\rho_t$ and $w_t$ are chosen, per row and per
cycle, to maximise the marginal likelihood of the new anomalies under the prior:
\begin{equation}
  (\rho_t,w_t)=\arg\max\ \log\mathcal N\!\Big(\y_t;\ \X_t\bb^{\rm prior},\ \sigma^2\big(\I+\X_t\mathbf P^{-1}\X_t^\top\big)\Big),
  \qquad \mathbf P = N\big(\rho\,\G^{\rm mem}_{t-1}+w\,\G^c\big),
  \label{eq:evidence}
\end{equation}
over a small grid ($\rho\in\{0.2,0.4,0.6,0.8,0.95\}$, $w/w_0\in\{0,0.1,0.3,1\}$), with
$\sigma^2$ profiled out. Each evaluation is one $N\times N$ Cholesky; the cost per cycle
is a small multiple of the plain EnKF-MC. The evidence penalises prior complexity
($\log\det$ term) and does not reward fitting noise, unlike cross-validation; it uses
only forecast anomalies, never observations or truth.

\paragraph{What the evidence cannot choose.} The tapering slope $\tau$ (and the base
ridge $\alpha$) are localization hyperparameters: they protect the \emph{analysis}
from noisy long-range coefficients, and the row likelihood does not see the analysis.
Chosen by the row evidence they go to zero and the filter degrades (measured,
Section~\ref{sec:results}). They are therefore treated like the radius and the
inflation: fixed by a sensitivity sweep.

\paragraph{Structure.} The predecessor set is the square of radius 2 (the optimum of
the uniform sweep at $N=40$) united with the climatological predecessors selected by a
lasso row regression on the $M$ climate states (window 6, universal-threshold penalty
with $c=0.5$). With the tapered ridge the long links cost nothing once the climatology
has been released; without it they do (Section~\ref{sec:results}).

\section{Testbed}\label{sec:setup}
Quasi-geostrophic model of Sakov and Oke as implemented in \texttt{pyteda}, $49\times49$
grid, dissipation ten times the default so that a stationary regime exists, state
$q$ only, Dirichlet boundary excluded from estimation ($47\times47=2209$ components),
$dt=0.3125$, 20 time units between analyses. Climatology: spin-up to $t=20000$ then
200 snapshots every 250 units (81 every 200 in the development runs); ensemble members
and truth are independent draws from it. Observations of $q$ (or of $\psi$, Sakov and
Oke's choice, through the Helmholtz inverse) with 5\% of the climatological spread as
error, on a fixed lattice of stride 2 (24\% of the interior) or on random networks
observing 1--40\% of it. Multiplicative inflation 1.05, calibrated for the plain
EnKF-MC; $\alpha=0.3$; $N=40$. Errors are RMSE of $q$ over the interior in units of
the climatological standard deviation, so that 1 is ``as bad as not assimilating''.

\section{Results}\label{sec:results}

\subsection{First analysis from a climatological ensemble}
At the first analysis the background error \emph{is} the climatology, and the
climatological prior mean is exact. Table~\ref{tab:first} gives the ratio of analysis to
background RMSE for five seeds, observing $\psi$, random network at 25\%.

\begin{table}[h]\centering\small
\begin{tabular}{@{}lrrrr@{}}\toprule
structure & ridge$\to0$ & ridge$\to\bb^c$, $\alpha=100$ \\\midrule
uniform $r=1$ & 0.707 & 0.691\\
uniform $r=2$ & 0.738 & 0.684\\
uniform $r=3$ & 0.760 & 0.669\\
uniform $r=4$ & 0.817 & \textbf{0.658}\\
climate lasso $w6$ & 0.709 & 0.660\\
LETKF $r=1$ & 0.792 & --\\\bottomrule
\end{tabular}
\caption{First analysis, $\psi$ observed, random network 25\%, $N=40$, five seeds.
With the climatological prior the ordering of the radii inverts: the wide row, useless
when estimated from 40 members, is the best when its coefficients come from 81 climate
states. The gain over the best plain row is 7\%. Across densities 1--40\% and both
observed variables the climatological prior is first at every density (2--9\%).}
\label{tab:first}
\end{table}

The dense climatological covariance of the same 81 states, used as an oracle $\B$ in a
Kalman analysis, gives 0.894 at 1\% density where the method gives 0.937 and the best
uniform row 0.971: the information limit of 22 observations is close, and the method
recovers about half of what separates the plain row from it.

\subsection{Cycled runs}
Figure~\ref{fig:rmse} shows 34 cycles on the lattice. The plain EnKF-MC with $r=2$
converges to 0.11; the LETKF, at its best radius, to 0.40; the Bayesian rows to 0.10,
below the plain rows at every cycle (ratio 0.80--0.95). Without the tapered ridge, the
climatological long links help at cycle 0 and then cost a plateau (0.71$\to$0.72$\to$0.72)
before recovering to 0.16; with it the descent is monotone from the best start.

\begin{figure}[h]\centering \fig{../results/figures_dev/bayes_rows_taper.png}{0.95}
\caption{Analysis RMSE per cycle, QG, $N=40$, lattice 24\%, inflation 1.05, one seed.}
\label{fig:rmse}\end{figure}

\subsection{What the evidence decides}
Figure~\ref{fig:maps} maps the hyperparameters chosen row by row. At cycle 0 the
climatological weight is at its maximum over the jet and the eddies and already
released in the quiescent north; from cycle 1 it is near zero everywhere but a few
rows. The forgetting factor is high (memory kept) in the north and low (forget) along
the jet. The implied background correlation of a point passes from the climatological
bands at cycle 0 to a local blob at cycle 15 with the same twelve neighbours in the
row: the values and the memory do the work, not the pattern.

\begin{figure}[h]\centering \fig{../results/figures_dev/bayes_rows_structure.png}{0.98}
\caption{Flow, climatological weight chosen at cycles 0 and 15, forgetting factor at
cycle 15, and the implied correlation of one point at cycles 0 and 15.}
\label{fig:maps}\end{figure}

\subsection{Sensitivity and what is fixed by hand}
$\rho$ and $w$ are chosen by the evidence; $w_0=M/N$ is a definition. Fixed: $\alpha=0.3$,
$r=2$, inflation 1.05 (inherited, sensitivity in the supplement), $\tau=0.3$
($\tau=1$ over-regularises: 0.20 at cycle 33; $\tau=0$ leaves the long links noisy),
and the climate lasso window and penalty (insensitive with the taper). Choosing
$\tau$ by the row evidence selects $\tau=0$ and the run ends at 0.25: the evidence of
a row is not the right criterion for a localization parameter.

\section{What does not work, and why}\label{sec:negative}
Five ways of choosing \emph{which} predecessors a row uses were tried on the same
testbed, each cycled for at least 36 cycles: a variational cost per component (which
selects the \emph{least} correlated candidate, a defect of its algebra), the partial
correlations of a wide probe precision, the trajectories of the ensemble-mean flow,
the lagged cross-correlation between analysis and forecast ensembles (an estimate of
the propagator), and a lasso row selection on the ensemble. All of them beat the
uniform square on a single analysis and none beats it after ten cycles. The reason is
visible in a perturbation experiment: a point perturbation of $q$ does not travel
along the flow within an assimilation interval, or within ten of them; it grows in
place. The converged forecast error is local, fast-decaying, and $r=2$ is the
Bickel--Levina optimum for it at $N=40$. Structures richer than that pay selection
noise without removing bias. The climatological structure is the one exception, and
only because the climatology, not the ensemble, supplies its coefficients.

\section{Discussion}\label{sec:discussion}
The method is a hybrid of climatological and ensemble information built inside the
precision, row by row, without a dense climatological covariance, without a blending
weight for covariances, and with the blending decided by the data where it can be. Its
cost is a small multiple of the EnKF-MC. Its gains are in the first analyses, in the
transient, and modestly in the converged state; they are measured on one seed and one
network and must be confirmed on sparse random networks with the LETKF alongside,
which is the experiment in progress. The evidence chooses memory and climatological
weight well; it cannot choose the tapering, and we do not claim it can.

\end{document}
